\chapter{V°Watch: values drawn in Rhino}\label{ch:watch}

\vn{Watch} is a list of expressions that are drawn in Rhino's viewports at every pause. It is a tab at
the bottom of the Script Editor and a page in its left panel, next to the editor's own Watch tab, which
shows values but doesn't draw them. Each editor has its own list, and it is kept between Rhino sessions.

\shot{01_watch_values_in_rhino}{Lesson 01 paused: \vn{Watch} lists \code{outline} (a PolylineCurve,
  \emph{drawn}) and \code{corners} (a list, \emph{drawn as points (4)}). Rhino's Top view draws them with
  their name dots: the outline crosses itself.}{fig:watch}

\section{Adding a watch}\label{sec:addwatch}

\begin{itemize}
  \item Type an expression in the box above the list and press \key{Enter}: a name (\code{crv}), an item
        (\code{curves[0]}), a slice (\code{pts[1:3]}), a member (\code{crv.PointAtStart}), a comprehension
        (\code{[c.PointAtEnd for c in curves]}) or a constructor (\code{rg.Line(a, b)}).
  \item Right-click a name in the code, or a row in the editor's Variables or Watch tab, and choose
        \ui{Watch in Rhino}. A nested row counts as its path, such as \code{curves[0]} or \code{node.x}.
  \item In a tooltip, click \ui{Watch in Rhino} (\see{sec:tooltiplinks}).
  \item In \vn{Python}, \code{watch('crv.PointAtStart')}.
  \item Right-click a member row inside a watch (\code{Degree}, \code{[0]}) and choose \ui{Watch in Rhino} to
        make it a watch of its own.
\end{itemize}

A watch added while paused is drawn at once. Each watch gets its own color.

\section{The list}\label{sec:watchlist}

The list is a tree, like the editor's Variables tab.

\begin{xltabular}{\linewidth}{@{}P{0.2\linewidth}L@{}}
\toprule
\ui{Show}       & An eye: open and blue means drawn, closed and gray means hidden. Click it to hide a watch without removing it. A hidden watch's row is grayed out. \\
\ui{Expression} & In its drawing color. Click the cell or press \key{F2} to edit it in place; the watch is drawn again with the new expression. \\
\ui{Value}, \ui{Type} & As the editor's Variables tab shows them. Rows expand into the value's members: a curve's \code{Degree}, \code{Domain}, \code{PointAtStart}; a list's \code{[0]}, \code{[1]}, … Rhino objects show their data as in the Variables tab. \\
\ui{In Rhino}   & What the last pause found: \code{drawn}, \code{drawn (3)} for three items, \code{drawn by show\_sensors (12)} for a visualizer, the value's type with \emph{nothing to draw}, or the error. \\
\bottomrule
\end{xltabular}

A row is grayed when the last pause couldn't evaluate it, and every row is grayed while the script runs
or after it ended (\emph{running}, \emph{not running}).

\paragraph{Toolbar.} One row above the list: the expression to add, then \ui{Add}, \ui{Remove} (or \key{Del}),
\ui{Bake}, \ui{Edit in Rhino}, \ui{Auto edit} (a switch with a green \ui{ON} or gray \ui{OFF} label) and
\ui{Clear all}. Their tooltips name them. The right-click menu has \ui{Interpret as}, \ui{In every thread}, \ui{Remove},
\ui{Bake to Rhino} and \ui{Edit in Rhino} (on a member row, \ui{Watch in Rhino}).

\section{What is drawn, and how}\label{sec:whatdrawn}

\begin{xltabular}{\linewidth}{@{}P{0.42\linewidth}L@{}}
\toprule
\textsf{\textbf{Value}} & \textsf{\textbf{Drawn as}} \\ \midrule
\code{Curve}, \code{Line}, \code{Polyline}, \code{Arc}, \code{Circle}, \code{Ellipse}, \code{Rectangle3d} & the curve \\
\code{Point3d}, \code{Point3f}, \code{Point} & a point \\
\code{Vector3d}, \code{Vector3f} & an arrow from the world origin \\
\code{Plane} & its axes, x red, y green, z blue, 60 pixels long at any zoom \\
\code{Box}, \code{BoundingBox}, \code{Brep}, \code{Surface}, \code{Extrusion} & shaded, translucent in the watch's color, with isocurves \\
\code{Mesh}, \code{SubD} & shaded, with mesh wires \\
\code{PointCloud}, \code{TextDot} & as such \\
a \code{Guid} of a Rhino object (rhinoscriptsyntax) & that object's geometry, in the watch's color \\
a .NET collection, a Python list, tuple or set & every drawable item, up to 4 levels deep and 10\,000 items \\
your own class & what its visualizer returns (\cref{ch:visualizers}) \\
raw data: numbers in tuples and lists & by \ui{Interpret as}: points, a polyline, vectors, … (\see{sec:interpretas}) \\
\bottomrule
\end{xltabular}

Other values aren't drawn; the \ui{In Rhino} column says what they are. A generator or other iterator
is left alone, since reading it would consume it.

\paragraph{Name dots.} Each drawn value has a dot with its name (and item count), up and to the right of
it, with a tail that curves down from the dot's bottom edge to a point at the value, like a speech
bubble. That way the dot never covers a point, at any zoom. Four settings in
\menu{V°Settings\>Rhino\>Watches in Rhino}: \setting{Name dots} (on or off; the values stay drawn),
\setting{Name dot offset} (pixels; 0 puts the dot on the value), \setting{Name dot opacity} (\%; lower
lets geometry show through) and \setting{Name dot pointer} (the curved \emph{Tail} or a straight \emph{Line}).

\section{Interpret as}\label{sec:interpretas}

Raw data, such as \code{[(0,0,0), (5,0,0), (5,5,0)]}, can mean several things. Right-click a watch and
choose \menu{Interpret as}: \emph{Auto} (the first, checked by default), or one of the choices that fit the
value, read from the paused script. A triple can be a point or a vector; a list of triples points, a
polyline, a closed polyline, an interpolated curve or vectors; pairs of triples lines or anchored
vectors; six-number rows anchored vectors; point triples triangles; containers are taken item by item.
\emph{Plain data} draws nothing. See \cref{ch:visualizers} for the full order.

\section{Selection, both ways}\label{sec:selection}

\begin{itemize}
  \item Selecting a row, in \vn{Watch} or in the editor's Variables and Watch tabs, draws its geometry in
        Rhino's selection color, and selects its objects being edited.
  \item Clicking drawn geometry in a viewport selects its row (a list item's row when its list is
        expanded) and its variable. \setting{Select a watch by clicking its drawing}: a plain click (the
        default), \keys{Alt+click}, or off.
  \item Selecting an object being edited in Rhino highlights its watch and selects its variable
        (\setting{Selecting edited objects highlights their watch}).
\end{itemize}

\section{Watches in every thread}

While a script runs more than one thread, each watch is also evaluated in the other threads and drawn
once per thread, as \code{pt @Worker-1}. The details are in \see{sec:everythread}.

\section{C\# and IronPython scripts}

For C\# scripts and IronPython~2 scripts, watches of plain variable names are drawn and baked;
expressions need Python~3.

\begin{lesson}{lessons 01, 06, 07}
\lessonfile{01_watch_geometry.py}: V°Watch. \lessonfile{06_watch_expressions.py}: expressions in watches.
\lessonfile{07_raw_data_interpret_as.py}: Interpret as.
\end{lesson}
